% Section 4. Controls. Dose detection and calibration details are in Supplement S5.

\section{Controls}
\label{sec:controls}

With every model failing every rung above the gate, the ladder has to be shown to leave real people unfailed and to catch failures of a known size.
The controls build pseudo-models out of real survey respondents.
Each replicate splits the NHANES adults at random into two halves within the 48 persona cells, uses one half as the reference, and fills the corpus design from the other, each draw a real respondent from the persona's cell, and four such pseudo-models stand in for the four language models.
Failures are then planted into the same respondents, each aimed at one rung, by making the items independent, shifting every persona mean, doubling the income gap, compressing answer patterns toward the typical one, and making three items independent within the low-income group.
Table~\ref{tab:panel} reports how often each rung passes for each simulator type.

% Written by scripts/88c_panel_table.py from analysis_brm/intermediate_panel/88_confusion.csv. Do not edit by hand.
\begin{table}[!ht]
\centering
\caption{Verdicts for Seven Simulator Types Built From Survey Respondents}
\label{tab:panel}
\small
\setlength{\tabcolsep}{4pt}
\begin{tabular}{lrrrrrrrrr}
\toprule
Simulator & Gate & L1 & L2 fail & L3 & R1 & R2 & R4 fail & R4 unres. & L4 \\
\midrule
Real respondents & 100.0 & 92.9 & 0.0 & 100.0 & 100.0 & 100.0 & 4.0 (25) & 96.0 & 0.0 \\ % R: analysis_brm/intermediate_panel/88_confusion.csv
Independent items & 100.0 & 0.0 & 0.0 & 100.0 & 0.0 & 0.0 & 0.0 (5) & 0.0 & 0.0 \\ % R: analysis_brm/intermediate_panel/88_confusion.csv
Shifted, +2.1 & 28.6 & 72.9 & 0.7 & 0.0 & 100.0 & 100.0 & 0.0 (5) & 100.0 & 0.0 \\ % R: analysis_brm/intermediate_panel/88_confusion.csv
Shifted, +4.7 & 0.0 & 15.7 & 0.7 & 0.0 & 100.0 & 100.0 & 0.0 (5) & 100.0 & 0.0 \\ % R: analysis_brm/intermediate_panel/88_confusion.csv
Income gap doubled & 100.0 & 84.3 & 92.5 & 85.4 & 100.0 & 100.0 & 0.0 (5) & 100.0 & 0.0 \\ % R: analysis_brm/intermediate_panel/88_confusion.csv
Compressed & 100.0 & 0.0 & 0.0 & 100.0 & 98.6 & 1.4 & 0.0 (5) & 100.0 & 0.0 \\ % R: analysis_brm/intermediate_panel/88_confusion.csv
Non-invariant (income) & 100.0 & 2.9 & 0.0 & 100.0 & 100.0 & 0.0 & 100.0 (25) & 0.0 & 0.0 \\ % R: analysis_brm/intermediate_panel/88_confusion.csv
\bottomrule
\end{tabular}
\begin{tabnote}
\textit{Note.} L1 to L4 are levels 1 to 4, and R1, R2 and R4 are the level-4 rules. Percentage of readings that pass, except the columns marked fail or unresolved. Gate, L2 and L3 are read per pseudo-model (280 readings per type); L1, R1 and R2 on one pseudo-model per replicate (70 replicates). R4 is read at margin .08 on the number of replicates in parentheses; it gives no verdict without a general factor. L2 passes in no reading of any type, because at 30 draws per persona the interval of even a faithful gap is wider than the kept region; its remaining readings are unresolved. L4 is the level-4 pass (R1, R2 and R4 in both framings). Shifted types raise every persona by the stated points; the doubled gap raises low-income personas only; non-invariant makes three items independent within the low-income group. Construction, seeds and per-step results: \texttt{paper\_brm/analysis\_brm/intermediate\_panel} in the release.
\end{tabnote}
\end{table}

\FloatBarrier

Real respondents pass the gate, level 3, R1 and R2 in every reading and level 1 in most, and level 2 never fails them, while each planted failure trips the rung it targets. % R: analysis_brm/intermediate_panel/88_confusion.csv
Their persona means reach a dependability of only .70 at 30 draws, since real demographic cells differ little, and the gate therefore requires $\phi \ge .80$ only of studies that compare or rank individual personas. % R: analysis_brm/intermediate_panel/88_panel_verdicts.csv
Two rungs also react to failures they do not target.
The gate fails most shifted samples, as more depressed respondents vary more between draws, and R2 fails compressed and non-invariant samples, both of which change the size of the loadings. % R: analysis_brm/intermediate_panel/88_confusion.csv
Each draw of a pseudo-model comes from a different respondent, and the controls therefore vary more between draws than a language model does.
The gate's reaction to a shifted sample need not carry over to models.
In a separate set of 200 replicates with graded doses, each rung detects a distortion smaller than the one the language models show at levels 1, 3 and 4, and the rate at which level 1 passes real respondents in those replicates sets the level-1 floor (Supplement S5).

The controls also show how much data each verdict takes.
A gap is kept only when the whole interval of its ratio lies inside the kept region, and an interval that narrow needs more draws and a more precise reference than one lying wholly outside the region.
On persona grids of up to 432 cells, level 2 rarely fails a faithful simulator built from real respondents and almost always fails a doubled income gap. % R: analysis_brm/panel_power/90_cells.csv; analysis_brm/panel_power/90b_summary.csv
At 30 draws per persona the income gap of a faithful simulator is still never read \vd{kept}, because draws of one persona vary enough to hold its interval wider than the kept region. % R: analysis_brm/panel_power/power_l2_contrast_readings.csv
With no simulation error the same gap is kept in nearly every replicate, and the Black--White and Hispanic--White gaps are never kept on any grid, because NHANES measures them too loosely. % R: analysis_brm/panel_power/power_l2_contrast_readings.csv
A researcher who needs a kept gap therefore plans for more draws per persona and for a population gap the survey estimates with a narrow interval, and a finer persona grid does not supply the second.
R4 shows the same asymmetry, leaving some invariance steps unresolved for real respondents in almost every replicate while failing every planted invariance failure, and a design that needs an R4 pass needs more personas per group than the corpus has. % R: analysis_brm/intermediate_panel/88_confusion.csv; analysis_brm/intermediate_panel/88_r4_steps.csv
